% 第8章 8.2.1：张量积网格。(a) 每轴结点 t_i=i/4（m=5），乘积帽形基 phi_1(x)phi_2(y) 的支撑 [t_0,t_2]x[t_1,t_3]（蓝底），
% 查询点 (0.6,0.35) 所在单元的四个角结点（绿）决定双线性插值；(b) 完整系数个数 m^d，m=4 与 m=10，d=1,...,6（对数纵轴）。
\begin{tikzpicture}[font=\figurefont\footnotesize]
\pgfplotsset{every axis/.style={axis lines=left,
  x axis line style={draw=FigureStroke,line width=.6pt},y axis line style={draw=FigureStroke,line width=.6pt},
  tick style={draw=FigureStroke},tick label style={font=\figurefont\scriptsize,text=FigureInk},
  title style={font=\figurefont\footnotesize},label style={font=\figurefont\scriptsize},
  legend style={draw=none,fill=none,font=\figurefont\scriptsize,cells={anchor=west}}}}
\begin{scope}[x=36mm,y=36mm]
\fill[FigureBlueFill] (0,0.25) rectangle (0.5,0.75);
\foreach \i in {0,...,4}{\draw[FigureGrid,line width=.5pt] (\i/4,0)--(\i/4,1); \draw[FigureGrid,line width=.5pt] (0,\i/4)--(1,\i/4);}
\draw[FigureStroke,line width=.6pt] (0,0) rectangle (1,1);
\foreach \i in {0,...,4}{\foreach \k in {0,...,4}{\fill[FigureMuted] (\i/4,\k/4) circle (0.9pt);}}
\fill[FigureBlue] (0.25,0.5) circle (1.6pt);
\node[anchor=south west,font=\figurefont\scriptsize,text=FigureBlue,inner sep=1pt] at (0.25,0.5) {$(t_1,t_2)$};
\draw[FigureGreen,line width=.9pt] (0.5,0.25) rectangle (0.75,0.5);
\foreach \p in {(0.5,0.25),(0.75,0.25),(0.5,0.5),(0.75,0.5)}{\fill[FigureGreen] \p circle (1.4pt);}
\draw[FigureInk,line width=.6pt] (0.6,0.35) circle (1.5pt);
\node[anchor=north,font=\figurefont\scriptsize] at (0.5,-0.04) {$x$};
\node[anchor=east,font=\figurefont\scriptsize] at (-0.04,0.5) {$y$};
\node[font=\figurefont\footnotesize] at (0.5,1.1) {(a) 乘积基与网格};
\end{scope}
\begin{semilogyaxis}[at={(54mm,0)},anchor=south west,width=50mm,height=40mm,xmin=0.6,xmax=6.4,ymin=1,ymax=3e6,
  xtick={1,...,6},xlabel={维数 $d$},ylabel={系数个数 $m^d$},title={(b) 组合规模},
  ymajorgrids,grid style={draw=FigureGrid,line width=.4pt},legend style={at={(0.03,0.97)},anchor=north west}]
\addplot[draw=FigureBlue,line width=1.1pt,mark=*,mark size=1.4pt,mark options={fill=FigureBlue,draw=FigureBlue,solid}] coordinates {(1,10) (2,100) (3,1000) (4,1e4) (5,1e5) (6,1e6)};
\addplot[draw=FigureGreen,line width=1.1pt,dashed,mark=square*,mark size=1.3pt,mark options={fill=FigureGreen,draw=FigureGreen,solid}] coordinates {(1,4) (2,16) (3,64) (4,256) (5,1024) (6,4096)};
\legend{$m=10$,$m=4$}
\end{semilogyaxis}
\end{tikzpicture}
